\uuid{qfI7}
\exo7id{4416}
\titre{exo7 4416}
\auteur{quercia}
\organisation{exo7}
\datecreate{2010-03-14}
\isIndication{false}
\isCorrection{true}
\chapitre{Série numérique}
\sousChapitre{Autre}
\cpge{mpsi-mp2i.procedes-sommatoires-discrets.a, pcsi.series-numeriques.a, ptsi.series-numeriques.a, mpsi-mp2i.procedes-sommatoires-discrets.b, pcsi.series-numeriques.b, ptsi.series-numeriques.b, mpsi-mp2i.procedes-sommatoires-discrets.c, pcsi.series-numeriques.c, ptsi.series-numeriques.c, mp-mpi.series-numeriques-vectorielles.b, pc.series-numeriques, psi.series-numeriques.a, pt.series-numeriques.a}
\module{Analyse}
\niveau{L1}
\difficulte{}

\contenu{
\question{
Soit $\alpha >0$. Étudier la série $\sum u_n$, avec 
$u_n=\frac{(-1)^n}{\sqrt{n^{\alpha}+(-1)^n}}\cdotp$
}
\reponse{
$u_n = \frac{(-1)^n}{n^{\alpha/2}}-\frac1{2n^{3\alpha/2}}
+  o\Bigl(\frac1{n^{3\alpha/2}}\Bigr)$,
il y a convergence ssi $\alpha>\frac23$.
}
}
